\section{Actual Haar admissibility and packet normalization}
\label{sec:iut-actual-haar-admissible-orbit}

The preceding source audit leaves a concrete local question: can the repaired
admissible-domain interface be realized by an actual measure, with the sign
and coefficient required by the source?  Proposition~3.9(i) of
\cite{IUT3} takes the nonarchimedean log-volume domain to be the nonempty
compact open subsets, normalizes the integral structure to log-volume zero,
and uses packet normalization so that multiplication by the local rational
prime subtracts its logarithm.  We now construct the underlying Haar model
and separate raw local volume from that packet normalization.

Let $K$ be a nontrivially normed field with proper metric, let $\mu$ be a
regular additive Haar measure, and let $a\in K^\times$.  The local fields
used below have this properness property.  Write
$m_a(x)=ax$ and let $\Delta_\mu:K^\times\to\R_{>0}$ be the distributive
Haar character, so that
\[
                 \mu(aU)=\Delta_\mu(a)\mu(U).
\]
Define the finite-positive domain
\[
 \mathcal F_\mu=
 \{U:U\text{ is measurable and }0<\mu(U)<\infty\},
 \qquad L_\mu(U)=\log\mu(U).
\]
The conversion from an extended nonnegative measure to a real logarithm is
made only after both inequalities in the definition of $\mathcal F_\mu$ have
been proved.

\begin{theorem}[Actual scalar-preimage law]
\label{thm:actual-haar-scalar-preimage}
For every $U\in\mathcal F_\mu$,
\[
 m_a^{-1}(U)=a^{-1}U\in\mathcal F_\mu,
 \qquad
 L_\mu(m_a^{-1}U)
 =L_\mu(U)+\log\Delta_\mu(a^{-1}).
 \tag{\ref{thm:actual-haar-scalar-preimage}.1}
\]
If $U$ is also compact and open, then $m_a^{-1}(U)$ is compact and open.
\end{theorem}

\begin{proof}
The carrier identity follows in both directions from
$ax\in U\Longleftrightarrow x=a^{-1}(ax)\in a^{-1}U$.  Multiplication by
$a$ is a homeomorphism, so its inverse image preserves measurability,
openness, and compactness.  The Haar change-of-variables identity gives
\[
 \mu(m_a^{-1}U)=\mu(a^{-1}U)
 =\Delta_\mu(a^{-1})\mu(U).
\]
The character is finite and strictly positive.  The right side is therefore
finite and positive, and taking real logarithms proves the formula.
\end{proof}

Suppose now that the valuation ring $\OO_K$ is a discrete valuation ring,
the residue field has cardinality $q$, and $\mu(\OO_K)=1$.  Let $\pi$ be a
uniformizer.

\begin{theorem}[Uniformizer orbit and its obstruction]
\label{thm:actual-haar-uniformizer-orbit}
Put $B_n=\pi^{-n}\OO_K$ for $n\ge0$.  Every $B_n$ is a nonempty compact
open member of $\mathcal F_\mu$, and
\[
 m_\pi^{-1}(B_n)=B_{n+1},
 \qquad L_\mu(B_n)=n\log q.                 \tag{\ref{thm:actual-haar-uniformizer-orbit}.1}
\]
Consequently $n\mapsto B_n$ is injective, and no single
$V\in\mathcal F_\mu$ contains every $B_n$.
\end{theorem}

\begin{proof}
The $q$ additive cosets of $\pi\OO_K$ partition $\OO_K$.  Translation
invariance and $\mu(\OO_K)=1$ imply
$\mu(\pi\OO_K)=q^{-1}$, hence
$\Delta_\mu(\pi)=q^{-1}$.  Theorem~\ref{thm:actual-haar-scalar-preimage}
therefore adds $\log q$ under a uniformizer preimage.  The unit ball is a
nonempty open ball; properness makes it compact, and scalar multiplication
preserves these properties.  Induction gives the displayed identities.
Since $q>1$, the log-volumes are strictly increasing, proving injectivity.
If every $B_n$ lay in one finite-positive $V$, monotonicity would give
$n\log q\le L_\mu(V)$ for every $n$, contrary to the Archimedean property.
\end{proof}

The last conclusion rules out only a bounded, finite, or periodic
implementation of the complete preimage orbit.  It does not conflict with
the least-hull construction for one fixed relatively compact region in
Remark~3.9.5(i) of \cite{IUT3}, and it does not rule out the Haar/IUT route.

The rational-prime coefficient requires one further normalization.  Let
$K/\Q_p$ have ramification index $e$ and residue degree $f$.  Choose a unit
$u$ with
\[
                  p=u\pi^e,\qquad q=p^f.
\]
Multiplication by $u$ preserves Haar measure.  Thus raw Haar log-volume
changes under rational-prime preimage by
\[
 \log\Delta_\mu(p^{-1})=e\log q=ef\log p.   \tag{\ref{thm:actual-haar-uniformizer-orbit}.2}
\]

\begin{theorem}[Local-degree and packet normalization]
\label{thm:iut-local-degree-packet-normalization}
Define
\[
                  L^{\rm pkt}_\mu(U)=\frac{L_\mu(U)}{ef}.
\]
Then
\[
 L^{\rm pkt}_\mu(p^{-1}U)=L^{\rm pkt}_\mu(U)+\log p.
 \tag{\ref{thm:iut-local-degree-packet-normalization}.1}
\]
For finitely many components, if $\sum_iw_i=1$, the weighted packet volume
$\sum_iw_iL^{\rm pkt}_{\mu_i}(U_i)$ also changes by exactly $\log p$ under
simultaneous rational-prime preimage.
\end{theorem}

\begin{proof}
Divide the raw identity $ef\log p$ by the positive local degree $ef$.
Every normalized component then has the same increment $\log p$, so the
weighted increment is $(\sum_iw_i)\log p=\log p$.
\end{proof}

This two-stage calculation explains which coefficient has to be realized in
the actual tensor components.  Ordinary weight normalization cannot perform
the first stage.

\begin{proposition}[Full-premise raw-volume counterexample]
\label{prop:iut-raw-weight-counterexample}
The condition $\sum_iw_i=1$ alone does not turn raw Haar log-volume into a
$\log p$ rational-prime shift.
\end{proposition}

\begin{proof}
Let $K/\Q_p$ be the unramified quadratic extension, take the one-element
packet, and give it weight one.  Then $e=1$, $f=2$, $q=p^2$, and all local
field, Haar, positivity, and weight premises hold.  The raw shift is
$\log(p^2)=2\log p\ne\log p$.  Division by $ef=2$ restores the conclusion
of Theorem~\ref{thm:iut-local-degree-packet-normalization}.
\end{proof}

Proposition~\ref{prop:iut-raw-weight-counterexample} retires exactly the
weaker raw-volume claim.  The corrected route remains active.  It must still
construct $p=u\pi^e$ and $q=p^f$ uniformly in every actual tensor component,
identify the component functional with $L_\mu/(ef)$, transport the
compact-open admissible class through the direct-sum and least-hull
operations, and then prove the horizontal same-pilot and Ind3 comparison.
None of these global requirements follows from a local change-of-variables
formula.
